% GENERATED FILE. Edit canonical lessons, then run npm run generate:books.
% canonical-source-hash: fnv1a64:002ab35a1ccc1ac8
% canonical-lessons: JA-C341-settoku-suru, JA-C341-unagasu, JA-C341-hikitomeru, JA-C341-narubeku, JA-C341-tatoeba

\chapter{Persuading Someone}
\label{ch:ja-persuading-someone}
\hlchaptermodality{341}

\begin{chapteropening}
\textbf{By the end of this chapter:} \emph{I can try to persuade someone and give my reasons.}

Complete the last lesson of chapter 341: I can try to persuade someone and give my reasons.
\end{chapteropening}

\section[\texorpdfstring{\ja{せっとくする} (settoku suru)}{せっとくする (settoku suru)}]{\ja{せっとくする} (settoku suru) — to persuade}

\label{lesson:JA-C341-settoku-suru}



\begin{quote}
Before the new one: say the Japanese for to make an effort, then the Japanese for a warning.
\end{quote}

\subsection*{You'll want to know: \ja{せっとくする}}
\textbf{\ja{せっとくする}} — \emph{settoku suru} — \textquotedblleft{}to persuade\textquotedblright{}.

\begin{quote}
\textbf{\ja{ははが} \ja{せっとくしたので、ちちは} \ja{やめました。}} — \emph{haha ga settoku shita node, chichi wa yamemashita} — \textquotedblleft{}My mother talked him round, so my father gave it up.\textquotedblright{}
\end{quote}

Polite \textbf{\ja{せっとくします}}, \emph{settokushimasu}; a \ja{する} verb. Past (polite): \textbf{\ja{ははが} \ja{ちちを} \ja{せっとくしました。}}, \emph{haha ga chichi o settoku shimashita}.

\begin{cousinweb}
\textbf{\ja{せっとくする}} (\emph{settoku suru}) is built from a Sino-Japanese noun and the native verb \emph{suru}, \textquotedblleft{}to do\textquotedblright{}. The noun, \emph{settoku}, is two Chinese characters read with their on'yomi, readings borrowed from Chinese: \emph{setsu}, \textquotedblleft{}explanation\textquotedblright{} (Middle Chinese \emph{sywet}) and \emph{toku}, \textquotedblleft{}gain\textquotedblright{} (\emph{tok}) (KANJIDIC2; JMdict; Middle Chinese in Baxter's transcription). The two join as \emph{set-toku}.
\end{cousinweb}

\begin{grammarlens}[title={What you've built}]
One more word for giving reasons and opinions.
\end{grammarlens}

\subsection*{Your turn}
\begin{itemize}
\raggedright
  \item \emph{Say it:} \emph{settoku suru}
  \item \emph{Say it:} \emph{settoku suru}, once more
  \item \emph{Say it:} \emph{haha ga settoku shita node, chichi wa yamemashita} — My mother talked him round, so my father gave it up.
  \item \emph{Recall:} say the Japanese for to make an effort, then the Japanese for a warning, then say \emph{settoku suru} again
\end{itemize}

\begin{tcolorbox}[breakable,colback=teal!4,colframe=teal!35!black,title={Before you move on}]
What does \ja{せっとくする} mean? (\textquotedblleft{}To persuade\textquotedblright{}.) Say it once more.
\end{tcolorbox}
\section[\texorpdfstring{\ja{うながす} (unagasu)}{うながす (unagasu)}]{\ja{うながす} (unagasu) — to urge; to prompt}

\label{lesson:JA-C341-unagasu}



\begin{quote}
Before the new one: say the Japanese for a warning, then the Japanese for to persuade.
\end{quote}

\subsection*{You'll want to know: \ja{うながす}}
\textbf{\ja{うながす}} — \emph{unagasu} — \textquotedblleft{}to urge; to prompt\textquotedblright{}.

\begin{quote}
\textbf{\ja{じかんが} \ja{ないので、せんせいが} \ja{がくせいを} \ja{うながしました。}} — \emph{jikan ga nai node, sensei ga gakusei o unagashimashita} — \textquotedblleft{}There was no time, so the teacher hurried the students along.\textquotedblright{}
\end{quote}

Polite \textbf{\ja{うながします}}, \emph{unagashimasu}; a godan verb. Past (polite): \textbf{\ja{せんせいが} \ja{うながしました。}}, \emph{sensei ga unagashimashita}.

\begin{cousinweb}
\textbf{\ja{うながす}} (\emph{unagasu}) is a native Japanese word (wago), inherited rather than borrowed from Chinese, and its older origin is not settled. Its kanji was attached to it later as a kun'yomi, a native reading given to a Chinese character (JMdict, KANJIDIC2).
\end{cousinweb}

\begin{grammarlens}[title={What you've built}]
One more word for giving reasons and opinions.
\end{grammarlens}

\subsection*{Your turn}
\begin{itemize}
\raggedright
  \item \emph{Say it:} \emph{unagasu}
  \item \emph{Say it:} \emph{unagasu}, once more
  \item \emph{Say it:} \emph{jikan ga nai node, sensei ga gakusei o unagashimashita} — There was no time, so the teacher hurried the students along.
  \item \emph{Recall:} say the Japanese for a warning, then the Japanese for to persuade, then say \emph{unagasu} again
\end{itemize}

\begin{tcolorbox}[breakable,colback=teal!4,colframe=teal!35!black,title={Before you move on}]
What does \ja{うながす} mean? (\textquotedblleft{}To urge; to prompt\textquotedblright{}.) Say it once more.
\end{tcolorbox}
\section[\texorpdfstring{\ja{ひきとめる} (hikitomeru)}{ひきとめる (hikitomeru)}]{\ja{ひきとめる} (hikitomeru) — to ask someone to stay}

\label{lesson:JA-C341-hikitomeru}



\begin{quote}
Before the new one: say the Japanese for to persuade, then the Japanese for to urge.
\end{quote}

\subsection*{You'll want to know: \ja{ひきとめる}}
\textbf{\ja{ひきとめる}} — \emph{hikitomeru} — \textquotedblleft{}to ask someone to stay\textquotedblright{}.

\begin{quote}
\textbf{\ja{あめが} \ja{ふったので、ともだちを} \ja{ひきとめました。}} — \emph{ame ga futta node, tomodachi o hikitomemashita} — \textquotedblleft{}It rained, so I asked my friend to stay.\textquotedblright{}
\end{quote}

Polite \textbf{\ja{ひきとめます}}, \emph{hikitomemasu}; an ichidan verb. Past (polite): \textbf{\ja{ははが} \ja{ともだちを} \ja{ひきとめました。}}, \emph{haha ga tomodachi o hikitomemashita}.

\begin{cousinweb}
\textbf{\ja{ひきとめる}} (\emph{hikitomeru}) is a native Japanese word (wago), inherited rather than borrowed from Chinese: it is built from \emph{hiki}, \textquotedblleft{}pulling (hiku)\textquotedblright{} and \emph{tomeru}, \textquotedblleft{}to stop\textquotedblright{}, and its kanji were attached later as kun'yomi (JMdict).
\end{cousinweb}

\begin{grammarlens}[title={What you've built}]
One more word for giving reasons and opinions.
\end{grammarlens}

\subsection*{Your turn}
\begin{itemize}
\raggedright
  \item \emph{Say it:} \emph{hikitomeru}
  \item \emph{Say it:} \emph{hikitomeru}, once more
  \item \emph{Say it:} \emph{ame ga futta node, tomodachi o hikitomemashita} — It rained, so I asked my friend to stay.
  \item \emph{Recall:} say the Japanese for to persuade, then the Japanese for to urge, then say \emph{hikitomeru} again
\end{itemize}

\begin{tcolorbox}[breakable,colback=teal!4,colframe=teal!35!black,title={Before you move on}]
What does \ja{ひきとめる} mean? (\textquotedblleft{}To ask someone to stay\textquotedblright{}.) Say it once more.
\end{tcolorbox}
\section[\texorpdfstring{\ja{なるべく} (narubeku)}{なるべく (narubeku)}]{\ja{なるべく} (narubeku) — as much as possible}

\label{lesson:JA-C341-narubeku}



\begin{quote}
Before the new one: say the Japanese for to urge, then the Japanese for to ask someone to stay.
\end{quote}

\subsection*{You'll want to know: \ja{なるべく}}
\textbf{\ja{なるべく}} — \emph{narubeku} — \textquotedblleft{}as much as possible\textquotedblright{}.

\begin{quote}
\textbf{\ja{あぶないから、なるべく} \ja{はやく} \ja{かえった} \ja{ほうが} \ja{いいです。}} — \emph{abunai kara, narubeku hayaku kaetta hō ga ii desu} — \textquotedblleft{}It's dangerous, so you had better go home as early as you can.\textquotedblright{}
\end{quote}

\begin{cousinweb}
\textbf{\ja{なるべく}} (\emph{narubeku}) is a native Japanese word (wago), inherited rather than borrowed from Chinese, and its older origin is not settled. JMdict notes that it is usually written in kana (JMdict).
\end{cousinweb}

\begin{grammarlens}[title={What you've built}]
One more word for giving reasons and opinions.
\end{grammarlens}

\subsection*{Your turn}
\begin{itemize}
\raggedright
  \item \emph{Say it:} \emph{narubeku}
  \item \emph{Say it:} \emph{narubeku}, once more
  \item \emph{Say it:} \emph{abunai kara, narubeku hayaku kaetta hō ga ii desu} — It's dangerous, so you had better go home as early as you can.
  \item \emph{Recall:} say the Japanese for to urge, then the Japanese for to ask someone to stay, then say \emph{narubeku} again
\end{itemize}

\begin{tcolorbox}[breakable,colback=teal!4,colframe=teal!35!black,title={Before you move on}]
What does \ja{なるべく} mean? (\textquotedblleft{}As much as possible\textquotedblright{}.) Say it once more.
\end{tcolorbox}
\section[\texorpdfstring{\ja{たとえば} (tatoeba)}{たとえば (tatoeba)}]{\ja{たとえば} (tatoeba) — for example}

\label{lesson:JA-C341-tatoeba}



\begin{quote}
Before the new one: say the Japanese for to ask someone to stay, then the Japanese for as much as possible.
\end{quote}

\subsection*{You'll want to know: \ja{たとえば}}
\textbf{\ja{たとえば}} — \emph{tatoeba} — \textquotedblleft{}for example\textquotedblright{}.

\begin{quote}
\textbf{\ja{たとえば、だいこんは} \ja{からだに} \ja{いいと} \ja{おもいます。}} — \emph{tatoeba, daikon wa karada ni ii to omoimasu} — \textquotedblleft{}For example, I think daikon is good for you.\textquotedblright{}
\end{quote}

\begin{cousinweb}
\textbf{\ja{たとえば}} (\emph{tatoeba}) is a native Japanese word (wago), inherited rather than borrowed from Chinese, and its older origin is not settled. Its kanji was attached to it later as a kun'yomi, a native reading given to a Chinese character (JMdict, KANJIDIC2).
\end{cousinweb}

\begin{grammarlens}[title={What you've built}]
One more word for giving reasons and opinions.
\end{grammarlens}

\subsection*{Your turn}
\begin{itemize}
\raggedright
  \item \emph{Say it:} \emph{tatoeba}
  \item \emph{Say it:} \emph{tatoeba}, once more
  \item \emph{Say it:} \emph{tatoeba, daikon wa karada ni ii to omoimasu} — For example, I think daikon is good for you.
  \item \emph{Recall:} say the Japanese for to ask someone to stay, then the Japanese for as much as possible, then say \emph{tatoeba} again
\end{itemize}

\begin{tcolorbox}[breakable,colback=teal!4,colframe=teal!35!black,title={Before you move on}]
What does \ja{たとえば} mean? (\textquotedblleft{}For example\textquotedblright{}.) Say it once more.
\end{tcolorbox}
